\SourceLeaf{363}%
circumscribed circle by \(AB\), and denotes a line \(=AC-AB\) by \(CF\). He then derives the following proportions by arguments in the form customary among ancient geometers.
\par
{\fontsize{10}{12}\selectfont
\[
\begin{aligned}
AC:AB&=5AB+2CF:3AB+CF=19AB+7CF:11AB+4CF\\
&=71AB+26CF:41AB+15CF=265AB+97CF:153AB+56CF\\
&=989AB+362CF:571AB+209CF
\end{aligned}
\]
}
From the next-to-last proportion it follows readily that \(AC:AB>265:153\). The last gives the following if one introduces a line \(BD=2AB-AC=AB-CF\).
\[
AC:AB=1351AB-362BD:780AB-209BD<1351:780
\]
We readily grant that Mr. M., who has made himself very familiar with the form in which ancient geometers presented arithmetical arguments, may indeed have guessed \textsc{Archimedes}'s line of thought. A decision on this is all the harder because such inquiries can be approached in a great many ways. There are also other indications that the great Greek must have known a number of truths and devices of a decidedly uncommon kind, even in higher arithmetic.

One question still remains. If \textsc{Archimedes} found his approximate values methodically. why did he go as far as the eleventh term for the larger ones, when he went only as far as the eighth for the smaller ones? One would think that for the larger ones he would have been content with the ninth term \(\frac{362}{209}\), which would in any case have been sufficient to establish the lower bound \(3\frac{10}{71}\). One might perhaps be tempted to conclude from this that \textsc{Archimedes} found the fraction \(\frac{1351}{780}\) after all by a kind of lucky accident, and that the simpler \(\frac{362}{209}\) escaped him. Mr. M. thinks that \textsc{Archimedes} chose the former fraction because it was the simplest with a numerator of the order of thousands, just as he chose the fraction \(\frac{265}{153}\) as the simplest of the order of hundreds. But this reason does not seem satisfactory to us. We find it more likely that he preferred the fraction \(\frac{1351}{780}\) because he found that it happened to make a convenient simplification possible later in the calculation. For the 24-gon, this meant that a very close bound for the ratio which, in our terms, is \(1:\operatorname{cotang}7^0\,30'\) could be expressed very %
